\documentclass[12pt, a4paper]{article}

% --- 常用宏包 ---
\usepackage[UTF8]{ctex}
\usepackage{amsmath, amsthm, amssymb, amsfonts}
\usepackage{geometry}
\usepackage{fancyhdr}
\usepackage{enumerate}
\usepackage{graphicx}
\usepackage{hyperref}
\usepackage{bm}
\usepackage{tcolorbox}

% --- 页面设置 ---
\geometry{left=2.5cm, right=2.5cm, top=2.5cm, bottom=2.5cm}
\pagestyle{fancy}
\fancyhf{}
\rhead{\today}
\lhead{近世代数作业}
\cfoot{\thepage}

% --- 环境定义 ---
\newtcolorbox{problem}[1]{
    colback=gray!5!white,
    colframe=gray!75!black,
    fonttitle=\bfseries,
    title=题目 #1
}

\newenvironment{solution}{
    \begin{proof}[\textbf{解}]
}{
    \end{proof}
}

% --- 个人信息 ---
\title{近世代数作业 03}
\author{李睿阳 (524120910059)}
\date{\today}

\begin{document}
\maketitle


\begin{problem}{1}
    将下面的置换写成不相交轮换的乘积，再表示为对换的乘积，给出置换的奇偶性。
    \begin{enumerate}[(i)]
        \item $\begin{pmatrix} 1 & 2 & 3 & 4 & 5 & 6 & 7 & 8 \\ 8 & 2 & 6 & 3 & 7 & 4 & 5 & 1 \end{pmatrix}$
        \item $\begin{pmatrix} 1 & 2 & 3 & 4 & 5 & 6 \\ 5 & 1 & 3 & 6 & 2 & 4 \end{pmatrix}$
    \end{enumerate}
\end{problem}
\begin{solution}
    \begin{enumerate}[(i)]
        \item $(1\ 8)(3\ 6\ 4)(5\ 7) = (1\ 8)(3\ 4)(3\ 6)(5\ 7)$。共有 4 个对换，是偶置换。
        \item $(1\ 5\ 2)(4\ 6) = (1\ 2)(1\ 5)(4\ 6)$。共有 3 个对换，是奇置换。
    \end{enumerate}
\end{solution}

\begin{problem}{2}
    对以下置换 $\sigma, \tau$，计算 $\tau\sigma\tau^{-1}$。
    $\sigma = (1, 3, 5, 6)$, $\tau = (2, 5, 4, 6)$。
\end{problem}
\begin{solution}
    $\tau = (2\ 5\ 4\ 6)$，则 $\tau^{-1} = (6\ 4\ 5\ 2)$。
    $\tau\sigma\tau^{-1} = (2\ 5\ 4\ 6)(1\ 3\ 5\ 6)(6\ 4\ 5\ 2) = (1\ 3\ 4\ 2)$。
\end{solution}

\begin{problem}{3}
    已知对任意 $\sigma \in S_n$ 及对换 $(i, j)$，若 $i, j$ 分属于 $\sigma$ 的两个不同轮换且长度均大于 $1$，则 $N((i, j) \circ \sigma) = N(\sigma) - 1$。
    
    证明：在下列情形下，该结论仍成立：
    \begin{enumerate}
        \item $i, j$ 分属于不同轮换，且其中一个为 $(i)$，另一个长度 $\ge 2$；
        \item $i, j$ 分属于两个单元素轮换 $(i)$ 和 $(j)$。
    \end{enumerate}
\end{problem}
\begin{solution}
    这里使用 $N(\sigma)$ 表示 $\sigma$ 分解为不相交轮换的个数（包含单元素轮换）。
    \begin{enumerate}
        \item 设 $i$ 所在的轮换为 $(i)$，$j$ 所在的轮换为 $(j, a_2, \dots, a_k)$ ($k \ge 2$)。
        则 $(i, j) \circ (i)(j, a_2, \dots, a_k) = (i, j, a_2, \dots, a_k)$。
        该操作将两个不相交轮换合并为一个，其余轮换不变，故 $N((i, j) \circ \sigma) = N(\sigma) - 1$。
        \item 设 $i, j$ 分属于两个单元素轮换 $(i)$ 和 $(j)$。
        则 $(i, j) \circ (i)(j) = (i, j)$。
        同样将两个不相交轮换合并为一个，故 $N((i, j) \circ \sigma) = N(\sigma) - 1$。
    \end{enumerate}
\end{solution}

\begin{problem}{4}
    设 $\mathbb{R}$ 是实数加群，定义映射 $\mathbb{R} \times \mathbb{R}^2 \to \mathbb{R}^2$：对任意 $r \in \mathbb{R}, (x, y) \in \mathbb{R}^2$，$r(x, y) = (x + ry, y)$。证明该映射是 $\mathbb{R}$ 对 $\mathbb{R}^2$ 的作用。
\end{problem}
\begin{solution}
    设 $G = (\mathbb{R}, +)$，$X = \mathbb{R}^2$。验证群作用的两个条件：
    \begin{enumerate}
        \item 恒等性：$G$ 的单位元为 $0$。对于任意 $(x, y) \in \mathbb{R}^2$，$0(x, y) = (x + 0 \cdot y, y) = (x, y)$。
        \item 结合性：对于任意 $r_1, r_2 \in \mathbb{R}$ 及 $(x, y) \in \mathbb{R}^2$：
        一方面，$(r_1 + r_2)(x, y) = (x + (r_1 + r_2)y, y)$；
        另一方面，$r_1(r_2(x, y)) = r_1(x + r_2y, y) = ((x + r_2y) + r_1y, y) = (x + (r_1 + r_2)y, y)$。
        因此 $(r_1 + r_2)(x, y) = r_1(r_2(x, y))$。
    \end{enumerate}
    综上，该映射是 $\mathbb{R}$ 在 $\mathbb{R}^2$ 上的群作用。
\end{solution}


\end{document}
